% title:          Third derivative from symmetric differences
% area:           derivatives, interpolation
% methods:        mean-value-theorem, rolle-theorem
% difficulty:     medium
% difficulty-ai:  easy
% status:
% review:         human
% --- history: all optional, leave EMPTY if unknown ---
% origin:         imc
% origin-date:    2005-07
% origin-number:  Day 2, Problem 4
% also-in:
% references:     Earliest known appearance (to the best of our knowledge): IMC 2005, Day 2, Problem 4 - https://www.imc-math.org.uk/imc2005/day2_solutions.pdf
% history:        ai
\begin{problem}
    Let \( f \in D^3\left(\mathbb{R},\mathbb{R}\right) \) (three times differentiable). Show that
    there exists \( c \in \left]-1;1\right[ \) with
    \[
        \frac{f'''\left(c\right)}{6} = \frac{f\left(1\right)-f\left(-1\right)}{2}-f'\left(0\right).
    \]
\end{problem}
\begin{solution}
    Let \( A := \frac{f\left(1\right)-f\left(-1\right)}{2}-f'\left(0\right) \in \mathbb{R} \). Now,
    we want to search for a polynomial \( P\left(x\right) \in \mathbb{R}\left[x\right] \) of degree
    at most \( 2 \) such that
    \[
        g_{P\left(x\right)} \colon \mathbb{R} \longrightarrow \mathbb{R}, \quad x \longmapsto
        f\left(x\right)-Ax^3+P\left(x\right)
    \]
    has a maximal number of zeros. Indeed, the problem asks us to search for a \( c \in
    \left]-1;1\right[ \) such that the third derivative of \( x \longmapsto f\left(x\right)-Ax^3 \)
    vanishes at \( c \). The idea is the following: to add more freedom and maximise the chance of
    having such a \( c \), we notice that adding a polynomial of degree at most \( 2 \) clearly does
    not change this third derivative, and we want to add such a polynomial so that the function has
    a maximal number of zeros, as each use of Rolle's theorem gives a zero of the derivative between
    two zeros of a function. So, by making computations, we wish to find a polynomial
    \( P\left(x\right) \) of degree at most \( 2 \) such that the function above has at least
    \( 4 \) zeros. Indeed, by applying Rolle's theorem three times, we get a zero of the third
    derivative, because we lose one zero each time we pass to the derivative. More precisely,
    \( 3 \) distinct zeros of the function also suffice, together with a zero of its derivative
    different from the two zeros of the derivative that Rolle's theorem gives between them: the
    derivative then has again \( 3 \) distinct zeros, and we conclude in the same way.
    \\\\
    Write \( P\left(x\right) = p_2x^2+p_1x+p_0 \) with \( p_0, p_1, p_2 \in \mathbb{R} \) (possibly
    \( p_2 = 0 \)). We want \( g_{P\left(x\right)} \) to vanish at \( -1 \), \( 0 \) and \( 1 \),
    and its derivative to vanish at \( 0 \) too. These are the points where the statement gives
    information, through \( f\left(-1\right) \), \( f\left(1\right) \) and \( f'\left(0\right) \),
    and the zero at \( 0 \) costs nothing, as the constant term \( p_0 \) takes care of it. They are
    also what Rolle's theorem needs to reach a point of \( \left]-1;1\right[ \): the three zeros
    \( -1 < 0 < 1 \) give two zeros of the derivative, one in \( \left]-1;0\right[ \) and one in
    \( \left]0;1\right[ \), and the zero of the derivative at \( 0 \) adds a third one between them,
    so that the second derivative has two zeros (again in \( \left]-1;0\right[ \) and
    \( \left]0;1\right[ \)) and hence the third derivative will have one zero in
    \( \left]-1;1\right[ \). As \( \left(g_{P\left(x\right)}\right)'\left(x\right) =
    f'\left(x\right)-3Ax^2+2p_2x+p_1 \) for all \( x \in \mathbb{R} \), these four conditions form
    the linear system
    \[
        \begin{cases}
            f\left(-1\right)+A+p_2-p_1+p_0 = 0, \\
            f\left(0\right)+p_0 = 0, \\
            f\left(1\right)-A+p_2+p_1+p_0 = 0, \\
            f'\left(0\right)+p_1 = 0
        \end{cases}
    \]
    in \( p_0 \), \( p_1 \) and \( p_2 \): four equations for three unknowns, which agree thanks to
    the value of \( A \) (the third equation minus the first gives \( p_1 =
    A-\frac{f\left(1\right)-f\left(-1\right)}{2} \), which is \( -f'\left(0\right) \) by the
    definition of \( A \), as the last equation asks). After doing the work, we arrive at
    \( p_2 = f\left(0\right)-\frac{f\left(1\right)+f\left(-1\right)}{2} \), \( p_1 =
    -f'\left(0\right) \) and \( p_0 = -f\left(0\right) \), i.e., the function
    \[
        g \colon \mathbb{R} \longrightarrow \mathbb{R}, \quad x \longmapsto
        f\left(x\right)-Ax^3+\left(f\left(0\right)-\left(\frac{f\left(1\right)+f\left(-1\right)}{2}\right)\right)x^2-f'\left(0\right)x-f\left(0\right)
    \]
    is in \( D^3\left(\mathbb{R},\mathbb{R}\right) \) (as the sum of \( f \) and a polynomial
    function) and satisfies (for the sake of the verification) \( g\left(0\right) = 0 \),
    \[
        g\left(1\right) =
        f\left(1\right)-A+f\left(0\right)-\left(\frac{f\left(1\right)+f\left(-1\right)}{2}\right)-f'\left(0\right)-f\left(0\right)
    \]
    \[
        =
        f\left(1\right)+f'\left(0\right)-\left(\frac{f\left(1\right)-f\left(-1\right)}{2}\right)-\left(\frac{f\left(1\right)+f\left(-1\right)}{2}\right)-f'\left(0\right)
        = 0
    \]
    and
    \[
        g\left(-1\right) =
        f\left(-1\right)+A+f\left(0\right)-\left(\frac{f\left(1\right)+f\left(-1\right)}{2}\right)+f'\left(0\right)-f\left(0\right)
    \]
    \[
        =
        f\left(-1\right)+\left(\frac{f\left(1\right)-f\left(-1\right)}{2}\right)-f'\left(0\right)-\left(\frac{f\left(1\right)+f\left(-1\right)}{2}\right)+f'\left(0\right)
        = 0,
    \]
    where, in both computations, we replaced \( A \) by its definition and cancelled the terms in
    \( f\left(0\right) \) for the second equality, then those in \( f'\left(0\right) \),
    \( f\left(1\right) \) and \( f\left(-1\right) \) for the last one. So, by applying twice Rolle's
    theorem to \( g \) on \( \left[-1;0\right] \) and \( \left[0;1\right] \) (\( g \) is
    differentiable), there exist \( c_1 \in \left]-1;0\right[ \) and \( c_2 \in \left]0;1\right[ \)
    with \( g'\left(c_1\right) = 0 = g'\left(c_2\right) \), where, for all \( x \in \mathbb{R} \),
    \[
        g'\left(x\right) =
        f'\left(x\right)-3Ax^2+\left(2f\left(0\right)-f\left(1\right)-f\left(-1\right)\right)x-f'\left(0\right).
    \]
    Notice (again for the verification) that \( g'\left(0\right) = f'\left(0\right)-f'\left(0\right)
    = 0 \), so, by again applying twice Rolle's theorem to \( g' \) on \( \left[c_1;0\right] \) and
    \( \left[0;c_2\right] \) (\( g' \) is differentiable, as \( g \in
    D^3\left(\mathbb{R},\mathbb{R}\right) \)), there exist \( c_3 \in \left]c_1;0\right[ \) and
    \( c_4 \in \left]0;c_2\right[ \) with \( g''\left(c_3\right) = 0 = g''\left(c_4\right) \). Apply
    Rolle's theorem again to \( g'' \) on \( \left[c_3;c_4\right] \) (\( g'' \) is differentiable,
    for the same reason): there exists \( c \in \left]c_3;c_4\right[ \subset \left]-1;1\right[ \)
    with \( g'''\left(c\right) = 0 \), i.e., as \( g'''\left(x\right) = f'''\left(x\right)-6A \) for
    all \( x \in \mathbb{R} \) (by differentiating the expression of \( g' \) above twice),
    \[
        \frac{f'''\left(c\right)}{6} = A =
        \frac{f\left(1\right)-f\left(-1\right)}{2}-f'\left(0\right),
    \]
    as desired.
\end{solution}
